\section{\tool Overview}
\label{sec:overview}


\begin{figure*}[tbh]
    \centering
    \includegraphics[width=0.99\linewidth]{figure/Overview.pdf}
    % \vspace{-1ex}
    \caption{Overview of \tool. Blue boxes are components of \tool.}
    \label{fig:tool-overview}
\end{figure*}

\tool operates by identifying causal relationships between faults in distributed systems and linking them to form potential cycles. It simulates 
chains of fault propagations 
in complex workloads quasi-statically by combining the causal relationships derived from fault injections on simpler workloads.

At the core of \tool is a counterfactual causality analysis. 
Counterfactual causality analysis involves conducting two experiments: 
one with a fault injected to observe its impact on the system execution (referred to as the \textit{interference}), 
and the other without the fault injected to observe the system's execution (referred to as the counterfactual -- i.e., what would the system execution be had the fault not been injected). 
Specifically, we compare the execution trace of a fault injection run with its corresponding profile run 
(i.e., running the same test without the injection) 
and identify any additional faults triggered, 
which are considered to have a counterfactual causal relationship 
with the injected fault.

\textbf{Workflow}. 
\Cref{fig:tool-overview} illustrates the workflow of \tool.
(\circledtext{1}) Given the Java bytecode of the target system, 
\tool's \textit{static analyzer} identifies 
system-specific fault injection points, 
and instruments fault injection and monitoring hooks. 
At runtime, 
\tool's \textit{runtime agent} 
performs fault injection and monitoring 
when a hook is encountered. 


During fault injection experiments (\circledtext{2}), 
a \textit{workload driver} selects combinations of faults and workloads 
(integration tests shipped with the target system) 
to execute both the fault injection run 
and the corresponding fault-free profile run (counterfactual experiment).
Note that, 
for each fault, 
\tool only uses the set of tests that 
can reach the fault during execution 
in fault injection runs. 


Next, the \textit{fault causality analysis} module (\circledtext{3}) compares the execution traces from the profile and injection runs 
to identify counterfactual causal relationships 
-- i.e., \textit{fault propagations}. 
These causal relationships are used by the 
\textit{three-phase allocation (3PA) protocol} 
to guide future workload and fault selection 
via feedback to the workload driver.

The discovered causal relationships are also forwarded to \tool's \textit{bug detector} (\circledtext{4}), 
which connects \textit{compatible} relationships and 
performs a \textit{beam search} for cycle detection 
to identify \vcs (\circledtext{5}).


In the following sections, 
we first explain 
(\Cref{sec:fia}) 
how to instrument faults 
and perform the core counterfactual causality analysis, 
and then (\Cref{sec:alloc}) how to efficiently allocate testing resources, 
and finally (\Cref{sec:vc-detection}) how to connect compatible causal relationships and detect cycles. 

